% PDF p. 70
\needspace{8\baselineskip}
\subsection{庞加莱圆盘模型}
\begin{statement}{定义 2.7}{庞加莱圆盘模型}
在几何学中，庞加莱圆盘模型也称共形圆盘模型，定义为 $\mathbb H^n$ 中的开单位圆盘
\begin{equation*}
  \mathbb H^n=\{x\in\R^n:\lVert x\rVert<1\}.
  \tag{2.46}\label{eq:2.46}
\end{equation*}
为了理解庞加莱圆盘，先介绍双曲几何中的一条性质。给定一条直线，在它之外取一点，过该点至少有两条直线与给定直线平行（图~\ref{fig:2.4}）。

图~\ref{fig:2.4} 用庞加莱圆盘直观展示了双曲平行公设。
\begin{itemize}
  \item \textbf{测地线} $L$ 是圆盘的竖直直径。
  \item \textbf{黑点} $p$ 在圆盘内且不在 $L$ 上。
  \item \textbf{曲线}是过 $p$ 的两条不同的双曲直线，由与圆盘边界正交的圆给出，均不与 $L$ 相交。
\end{itemize}
% PDF p. 71
\begin{figure}[htbp]
  \centering
  \begin{tikzpicture}[x=3.25cm,y=3.25cm]
  \draw[thick] (0,0) circle (1);
  \draw[red!85!black,thick] (0,-1)--(0,1);
  \pgfmathsetmacro{\arcx}{1.14}
  \pgfmathsetmacro{\arcy}{.60}
  \pgfmathsetmacro{\arcr}{sqrt(\arcx*\arcx+\arcy*\arcy-1)}
  \pgfmathsetmacro{\middleangle}{atan2(\arcy,\arcx)+180}
  \pgfmathsetmacro{\anglewidth}{asin(1/sqrt(\arcx*\arcx+\arcy*\arcy))}
  \pgfmathsetmacro{\startangle}{\middleangle-\anglewidth}
  \pgfmathsetmacro{\endangle}{\middleangle+\anglewidth}
  \foreach \sign in {1,-1} {
    \begin{scope}[yscale=\sign]
      \draw[blue!65!black,thick]
      ({\arcx+\arcr*cos(\startangle)},{\arcy+\arcr*sin(\startangle)})
      arc[start angle=\startangle,end angle=\endangle,radius=\arcr];
    \end{scope}
  }
  \fill ({\arcx-sqrt(\arcr*\arcr-\arcy*\arcy)},0) circle (1pt)
    node[above right,inner sep=2pt] {$p$};
  \node[draw=black!25,inner sep=2pt,anchor=south west,font=\small] at (-1,1)
    {\tikz[baseline=-.35ex]\draw[red!85!black,thick] (0,0)--(.24cm,0); 直线 $L$};
\end{tikzpicture}

  \FigureCaption{2.4}{庞加莱圆盘。过 $p$ 的两条双曲平行线。}
\end{figure}

因此，过同一点 $p$ 至少存在两条不与 $L$ 相交的直线。这是双曲几何的一个特征。

在双曲几何中，直线表示为端点处与圆盘边界垂直的圆弧，直径也是允许的。类似于欧氏几何中的平行线，互不相交的两条圆弧对应于平行射线。正交相交的两条圆弧对应于垂线（见 Wolfram MathWorld 的 Poincaré Hyperbolic Disk）。

双曲几何中没有矩形。双曲几何的另一条重要性质是，任一三角形的内角和小于 $180^\circ$。
\end{statement}

\begin{statement}{命题 2.2}{}
庞加莱圆盘模型配有如下黎曼度量。
\begin{equation*}
  g=\frac{4\,dx\cdot dx}{(1-\lVert x\rVert^2)^2}.
  \tag{2.47}\label{eq:2.47}
\end{equation*}

\textbf{证明。}庞加莱圆盘模型是双曲几何的一种表示，其中双曲空间 $\mathbb H^2$ 映到 $\R^2$ 中的单位圆盘 $D^2$。回顾双曲空间 $\mathbb H^2$ 的定义为
\begin{equation*}
  -x_0^2+x_1^2+x_2^2=-1,\qquad x_0>0,
  \tag{2.48}\label{eq:2.48}
\end{equation*}
% PDF p. 72
从而
\begin{equation*}
  x_1^2+x_2^2=x_0^2-1
  \quad\text{或}\quad
  \frac{x_1^2}{x_0^2-1}+\frac{x_2^2}{x_0^2-1}=1.
  \tag{2.49}\label{eq:2.49}
\end{equation*}
注意
\[
  x_0+1>\sqrt{(x_0+1)(x_0-1)}=\sqrt{x_0^2-1}.
\]
因而
\begin{equation*}
  \frac{1}{(x_0+1)^2}<\frac{1}{x_0^2-1}.
  \tag{2.50}\label{eq:2.50}
\end{equation*}
由式~\eqref{eq:2.49} 与式~\eqref{eq:2.50} 得
\begin{equation*}
  \frac{x_1^2}{(x_0+1)^2}+\frac{x_2^2}{(x_0+1)^2}
  <\frac{x_1^2}{x_0^2-1}+\frac{x_2^2}{x_0^2-1}=1.
  \tag{2.51}\label{eq:2.51}
\end{equation*}
令
\begin{equation*}
  u=\frac{x_1}{1+x_0}\quad\text{且}\quad v=\frac{x_2}{1+x_0}.
  \tag{2.52}\label{eq:2.52}
\end{equation*}
结合式~\eqref{eq:2.51} 和式~\eqref{eq:2.52}，得到
\begin{equation*}
  u^2+v^2<1.
  \tag{2.53}\label{eq:2.53}
\end{equation*}
用映射\translatornote{原书在本段用 $\varnothing$ 记映射名称。式~\eqref{eq:2.49} 右式和式~\eqref{eq:2.50} 的分母在 $x_0=1$ 时为零。此时 $(x_0,x_1,x_2)=(1,0,0)$ 直接映到圆盘的原点。}
\begin{equation*}
  \varnothing\colon(x_0,x_1,x_2)\longmapsto
  \left(\frac{x_1}{1+x_0},\frac{x_2}{1+x_0}\right),
  \tag{2.54}\label{eq:2.54}
\end{equation*}
得到庞加莱圆盘
\begin{equation*}
  D^2=\{(u,v)\in\R^2:u^2+v^2<1\}.
  \tag{2.55}\label{eq:2.55}
\end{equation*}

接下来说明如何利用映射 $\varnothing$ 的拉回和双曲面度量，在圆盘上诱导黎曼度量。\translatornote{映射 $\varnothing$ 从双曲面指向圆盘。将双曲面度量拉回到圆盘时，用的是 $\varnothing^{-1}$。原文随后先计算圆盘欧氏内积经 $\varnothing$ 拉回到双曲面所得的双线性形式。}

回顾在微分几何中，映射 $\varnothing$ 的拉回可把微分形式、度量等张量从一个流形拉回到另一个流形。拉回记为 $\varnothing^*$，作用在微分形式或张量的空间上。

设 $\varnothing\colon M\to N$ 是两个流形之间的光滑映射，$g$ 是 $N$ 上的度量。$N$ 上微分形式（或度量）$g$ 的\textbf{拉回} $\varnothing^*$ 定义为
\begin{equation*}
  \varnothing^*g(\rho,w)=g\bigl(d\varnothing(\rho),d\varnothing(w)\bigr),
  \tag{2.56}\label{eq:2.56}
\end{equation*}
% PDF p. 73
其中 $\rho,w\in T_pM$ 是点 $p\in M$ 处的切向量，$d\varnothing\colon T_pM\to T_{\varnothing(p)}N$ 是映射 $\varnothing$ 的推前。

回顾
\[
  \begin{pmatrix}u\\v\end{pmatrix}
  =\varnothing(x_0,x_1,x_2)
  =\begin{pmatrix}\varnothing_1(x_0,x_1,x_2)\\
                   \varnothing_2(x_0,x_1,x_2)\end{pmatrix}
  =\begin{pmatrix}\dfrac{x_1}{1+x_0}\\[5pt]\dfrac{x_2}{1+x_0}\end{pmatrix}.
\]
现在计算雅可比矩阵 $d\varnothing$。
\begin{equation*}
  \begin{aligned}
    d\varnothing
      &=\begin{pmatrix}du\\dv\end{pmatrix}
      =\begin{pmatrix}
        \dfrac{\partial u}{\partial x_0}\,dx_0+\dfrac{\partial u}{\partial x_1}\,dx_1\\[5pt]
        \dfrac{\partial v}{\partial x_0}\,dx_0+\dfrac{\partial v}{\partial x_2}\,dx_2
       \end{pmatrix}\\[4pt]
      &=\begin{pmatrix}
        -\dfrac{x_1\,dx_0}{(1+x_0)^2}+\dfrac{dx_1}{1+x_0}\\[5pt]
        -\dfrac{x_2\,dx_0}{(1+x_0)^2}+\dfrac{dx_2}{1+x_0}
       \end{pmatrix}.
  \end{aligned}
  \tag{2.57}\label{eq:2.57}
\end{equation*}

设 $\rho=(\rho_0,\rho_1,\rho_2)\in T_pM$，$w=(w_0,w_1,w_2)\in T_pM$。于是
\begin{equation*}
  \begin{aligned}
    \rho_0&=dx_0,&\rho_1&=dx_1,&\rho_2&=dx_2,\\
    w_0&=dx_0,&w_1&=dx_1,&w_2&=dx_2.
  \end{aligned}
  \tag{2.58}\label{eq:2.58}
\end{equation*}
把式~\eqref{eq:2.58} 代入式~\eqref{eq:2.57}，得到
\begin{equation*}
  d\varnothing(\rho)=\begin{pmatrix}du\\dv\end{pmatrix}
  =\begin{pmatrix}
    -\dfrac{x_1\rho_0}{(1+x_0)^2}+\dfrac{\rho_1}{1+x_0}\\[5pt]
    -\dfrac{x_2\rho_0}{(1+x_0)^2}+\dfrac{\rho_2}{1+x_0}
  \end{pmatrix},
  \tag{2.59}\label{eq:2.59}
\end{equation*}
\begin{equation*}
  d\varnothing(w)
  =\begin{pmatrix}
    -\dfrac{x_1w_0}{(1+x_0)^2}+\dfrac{w_1}{1+x_0}\\[5pt]
    -\dfrac{x_2w_0}{(1+x_0)^2}+\dfrac{w_2}{1+x_0}
  \end{pmatrix}.
  \tag{2.60}\label{eq:2.60}
\end{equation*}

回顾庞加莱圆盘中的点 $(u,v)$ 属于 $\R^2$。因而，黎曼度量即内积，取 $d\varnothing(\rho)$ 与 $d\varnothing(w)$ 之间的度量。$g(d\varnothing(\rho),d\varnothing(w))$ 由下式给出。
% PDF p. 74
\begin{equation*}
  \begin{aligned}
    &g\bigl(d\varnothing(\rho),d\varnothing(w)\bigr)\\
    &\quad=\left(-\frac{x_1\rho_0}{(1+x_0)^2}+\frac{\rho_1}{1+x_0}\right)
      \left(-\frac{x_1w_0}{(1+x_0)^2}+\frac{w_1}{1+x_0}\right)\\
    &\qquad+\left(-\frac{x_2\rho_0}{(1+x_0)^2}+\frac{\rho_2}{1+x_0}\right)
      \left(-\frac{x_2w_0}{(1+x_0)^2}+\frac{w_2}{1+x_0}\right).
  \end{aligned}
  \tag{2.61}\label{eq:2.61}
\end{equation*}
双曲面 $\mathbb H^2$ 上洛伦兹度量所对应的拉回度量 $g^*(\rho,w)$ 为
\begin{equation*}
  g^*(\rho,w)=g\bigl(d\varnothing(\rho),d\varnothing(w)\bigr).
  \tag{2.62}\label{eq:2.62}
\end{equation*}
最后化简式~\eqref{eq:2.61}，得到 $g^*(\rho,w)$。展开式~\eqref{eq:2.61} 的第一项，得
\begin{equation*}
  \begin{aligned}
    &\left(-\frac{x_1\rho_0}{(1+x_0)^2}+\frac{\rho_1}{1+x_0}\right)
     \left(-\frac{x_1w_0}{(1+x_0)^2}+\frac{w_1}{1+x_0}\right)\\
    &\quad=\frac{x_1^2\rho_0w_0}{(1+x_0)^4}
      -\frac{x_1\rho_0w_1}{(1+x_0)^3}
      -\frac{x_1\rho_1w_0}{(1+x_0)^3}
      +\frac{\rho_1w_1}{(1+x_0)^2}.
  \end{aligned}
  \tag{2.63}\label{eq:2.63}
\end{equation*}
类似地，展开式~\eqref{eq:2.62} 的第二项。\translatornote{原文引用式~\eqref{eq:2.62}。此处展开的是式~\eqref{eq:2.61} 右侧的第二项。}
\begin{equation*}
  \begin{aligned}
    &\left(-\frac{x_2\rho_0}{(1+x_0)^2}+\frac{\rho_2}{1+x_0}\right)
     \left(-\frac{x_2w_0}{(1+x_0)^2}+\frac{w_2}{1+x_0}\right)\\
    &\quad=\frac{x_2^2\rho_0w_0}{(1+x_0)^4}
      -\frac{x_2\rho_0w_2}{(1+x_0)^3}
      -\frac{x_2w_0\rho_2}{(1+x_0)^3}
      +\frac{\rho_2w_2}{(1+x_0)^2}.
  \end{aligned}
  \tag{2.64}\label{eq:2.64}
\end{equation*}
将式~\eqref{eq:2.63} 与式~\eqref{eq:2.64} 相加，得到
\begin{equation*}
  \begin{aligned}
    &\frac{(x_1^2+x_2^2)\rho_0w_0}{(1+x_0)^4}
       -\frac{\rho_0(x_1w_1+x_2w_2)}{(1+x_0)^3}\\
    &\qquad-\frac{w_0(x_1\rho_1+x_2\rho_2)}{(1+x_0)^3}
       +\frac{\rho_1w_1+\rho_2w_2}{(1+x_0)^2}.
  \end{aligned}
  \tag{2.65}\label{eq:2.65}
\end{equation*}
由式~\eqref{eq:2.52} 得
\begin{equation*}
  \begin{aligned}
    1-u^2-v^2&=1-\left(\frac{x_1}{1+x_0}\right)^2
                 -\left(\frac{x_2}{1+x_0}\right)^2
      =\frac{(1+x_0)^2-x_1^2-x_2^2}{(1+x_0)^2}\\
    &=\frac{(1+x_0)^2-x_0^2+1}{(1+x_0)^2}
      =\frac{2(1+x_0)}{(1+x_0)^2}
      =\frac{2}{1+x_0}.
  \end{aligned}
  \tag{2.66}\label{eq:2.66}
\end{equation*}
由式~\eqref{eq:2.66} 可知
\begin{equation*}
  (1-u^2-v^2)^2=\frac{4}{(1+x_0)^2}.
  \tag{2.67}\label{eq:2.67}
\end{equation*}

% PDF p. 75
于是
\begin{equation*}
  \frac{4}{(1-u^2-v^2)^2}=(1+x_0)^2.
  \tag{2.68}\label{eq:2.68}
\end{equation*}
可以证明
\begin{equation*}
  du^2+dv^2=\frac{1}{(1+x_0)^2}\varnothing^*g(\rho,w).
  \tag{2.69}\label{eq:2.69}
\end{equation*}
$D^2$ 上的黎曼度量为
\begin{equation*}
  g=\frac{4}{(1-u^2-v^2)^2}(du^2+dv^2).
  \tag{2.70}\label{eq:2.70}
\end{equation*}
\end{statement}
